\subsection{A CRITERION FOR SOLVABILITY}

Lemma 2. Let \(x\) be an endomorphism of a finite-dimensional vector space \(V\) and \(s\) (resp. \(n\) ) its semi-simple (resp. nilpotent) component (cf.~Algebra, Chapter VIII, § 9, no. 4, Definition 4). Let \(ad_x\) , \(ad_s\) , \(ad_n\) be the respective images of \(x\) , \(s\) , \(n\) in the adjoint representation of \(gl(V)\) . Then \(ad_s\) (resp. \(ad_n\) ) is the semi-simple (resp. nilpotent) component of \(ad_x\) and is equal to a polynomial in \(ad_x\) with coefficients in \(K\) and no constant term.

We know that \(\operatorname{ad} x = \operatorname{ad} s + \operatorname{ad} n\) , \([\operatorname{ad} s, \operatorname{ad} n] = 0\) and \(\operatorname{ad} n\) is nilpotent (§ 4, Lemma 1). We show that \(\operatorname{ad} s\) is semi-simple. It suffices to do this when \(\mathbf{K}\) is algebraically closed (cf.~Algebra, Chapter VIII, § 9, no. 2, Proposition 3). Then let \((e_i)_{1 \leq i \leq n}\) be a basis of \(V\) such that \(s(e_i) = \lambda_i e_i\) ( \(\lambda_i \in K\) ). Let \((E_{ij})\) be the canonical basis of \(\mathbf{M}_n(K) = \mathfrak{gl}(V)\) . By formulae (5) of § 1,

\[
(\mathrm{ad} s). E _ {i j} = (\lambda_ {i} - \lambda_ {j}) E _ {i j}
\]

and hence ad s is semi-simple. The last assertion of the lemma follows from Algebra, Chapter VIII, § 9, no. 4, Proposition 8.

Lemma 3. Let \(\mathbf{M}\) be a finite-dimensional vector space, \(\mathbf{A}\) and \(\mathbf{B}\) two vector subspaces of \(\mathfrak{gl}(\mathbf{M})\) such that \(\mathbf{B} \subset \mathbf{A}\) and \(\mathbf{T}\) the set of \(t \in \mathfrak{gl}(\mathbf{M})\) such that \([t, \mathbf{A}] \subset \mathbf{B}\) . If \(z \in \mathbf{T}\) is such that \(\operatorname{Tr}(zu) = 0\) for all \(u \in \mathbf{T}\) , then \(z\) is nilpotent.

It suffices to prove this when \(\mathbf{K}\) is algebraically closed, which we shall assume henceforth. Let \(s\) and \(n\) be the semi-simple and nilpotent components of \(z\) and let \((e_i)\) be a basis of \(\mathbf{M}\) such that \(s(e_i) = \lambda_i e_i\) ( \(\lambda_i \in \mathbf{K}\) ). Let \(\mathbf{V} \subset \mathbf{K}\) be the vector space over \(\mathbf{Q}\) generated by the \(\lambda_i\) . We need to prove that \(\mathbf{V} = \{0\}\) . Let \(f\) be a \(\mathbf{Q}\) -linear form on \(\mathbf{V}\) and let \(t\) be the endomorphism of \(\mathbf{M}\) such that \(te_{i}=f(\lambda_{i})e_{i}\) . If \((E_{ij})\) is the canonical basis of gl(M) defined by \(E_{ij}e_{k}=\delta_{jk}e_{i}\) , then

\[
\begin{array}{l} (\mathrm{ad} s) E _ {i j} = (\lambda_ {i} - \lambda_ {j}) E _ {i j} \\ (\mathrm{ad} t) E _ {i j} = (f (\lambda_ {i}) - f (\lambda_ {j})) E _ {i j}. \end{array}
\]

There exists a polynomial \(\mathbf{P}\) with no constant term and with coefficients in \(\mathbf{K}\) such that \(\mathrm{P}(\lambda_i - \lambda_j) = f(\lambda_i) - f(\lambda_j)\) for all \(i\) and \(j\) (for if \(\lambda_i - \lambda_j = \lambda_h - \lambda_k\) , then \(f(\lambda_i) - f(\lambda_j) = f(\lambda_h) - f(\lambda_k)\) ) and, if \(\lambda_i - \lambda_j = 0\) , \(f(\lambda_i) - f(\lambda_j) = 0\) ). Then ad \(t = \mathrm{P}(\mathrm{ad}s)\) . On the other hand, ad \(s\) is a polynomial with no constant term in ad \(z\) . Now (ad \(z\) ) (A) \(\subset \mathrm{B}\) , whence also (ad \(t\) ) (A) \(\subset \mathrm{B}\) . By the hypothesis \(0 = \operatorname{Tr}(zt) = \sum_{\lambda_i} f(\lambda_i)\) , whence \(0 = f(\operatorname{Tr}(zt)) = \sum f(\lambda_i)^2\) . Since the \(f(\lambda_i)\) are rational numbers, \(f = 0\) , which completes the proof.

THEOREM 2 (Cartan's criterion). Let \(\mathfrak{g}\) be a Lie algebra, \(\mathbf{M}\) a finite-dimensional vector space, \(\rho\) a representation of \(\mathfrak{g}\) on \(\mathbf{M}\) and \(\beta\) the bilinear form on \(\mathfrak{g}\) associated with \(\rho\) . Then \(\rho(\mathfrak{g})\) is solvable if and only if \(\mathcal{D}\mathfrak{g}\) is orthogonal to \(\mathfrak{g}\) with respect to \(\beta\) .

It can obviously be reduced to the case where \(\mathfrak{g}\) is a Lie subalgebra of \(\mathfrak{gl}(\mathbf{M})\) and \(\rho\) is the identity mapping. If \(\mathfrak{g}\) is solvable, \(\mathcal{D}\mathfrak{g}\) is contained in the largest nilpotency ideal of the identity representation of \(\mathfrak{g}\) (Theorem 1) and hence is orthogonal to \(\mathfrak{g}\) with respect to \(\beta\) ( \(\S 4\) , Proposition 4 (d)). Suppose that \(\mathcal{D}\mathfrak{g}\) is orthogonal to \(\mathfrak{g}\) with respect to \(\beta\) . We prove that \(\mathfrak{g}\) is solvable. Let \(\mathrm{T}\) be the set of \(t\in \mathfrak{gl}(\mathbf{M})\) such that \([t,\mathfrak{g}]\subset \mathcal{D}\mathfrak{g}\) . If \(t\in \mathrm{T}\) and \(x,y\) belong to \(\mathfrak{g}\) , then \([t,x]\in \mathcal{D}\mathfrak{g}\) and hence

\[
\operatorname{Tr} (t [ x, y ]) = \beta ([ t, x ], y) = 0
\]

whence by linearity \(\mathrm{Tr}(tu) = 0\) for all \(u\in \mathcal{D}\mathfrak{g}\) . Also, clearly \(\mathcal{D}\mathfrak{g}\subset T\) . Hence (Lemma 3) every element of \(\mathcal{D}\mathfrak{g}\) is nilpotent. It follows that \(\mathcal{D}\mathfrak{g}\) is nilpotent (§ 4, Corollary 3 to Theorem 1) and hence that \(\mathfrak{g}\) is solvable (no. 3, Corollary 5 to Theorem 1).

